%=====================================================================
\section{Conclusion}
\label{sec:conclusion}

Decision orchestration separates deterministic dispatch and assembly from local model judgments. In the reported Chinese UMR runs, supervised artifact updates, task-level threshold tuning and decision routing improve several configurations without weight updates. Traces support inspection, not causal attribution. Training gains exceed test gains; matched baselines, repeated trajectories and targeted ablations remain necessary to establish the robustness and contribution of each component.

\section*{Limitations}
Our main evidence comes from one task in one language. The second benchmark (\S\ref{sec:res-second}) has not been run yet, and the English \umr runs, made with earlier versions of the procedure, are not reported. Test sets are small (99--157 sentences) and differ between stages of the work. Unifying them is planned (\exper{E1}). Most configurations were run once, and run-to-run noise (up to 0.6 \F for closed mode and 2.3 \F for a 9B model in fragment mode) is of the same order as some of the differences we report. The seeds were written by the authors, so their quality, and the effort it took to write them, affect every result. We have not measured that effort. Several of the models we use, including all three decision-model services and the frontier LLMs, are commercial APIs whose behavior may change; the release must include shareable predictions, exact versions and permitted cached answers before these runs can be called independently reproducible \todo{check that the release includes them}. The routing results come from 20 sentences, and the ``human'' in routing will be simulated from gold, which overstates real annotators' accuracy. Finally, AnCast and our fine-grained judge depend on an automatic alignment, and the gold itself contains errors. One gold graph has unbalanced brackets and is skipped by the scorer.

\section*{Ethics Statement}
We use the Chinese \umr annotations supplied in the research checkout. \todo{Verify the exact release, license, cleaned-gold provenance, redistribution and commercial-API use permissions.} We collected no new data from people and no personal information beyond what the released documents contain. Our goal is to reduce the cost of producing the first annotations for a new scheme, which could reduce demand for annotation work. We see the method as a way to direct annotators' effort to the decisions that need human judgment. Running the experiments in this paper used commercial APIs, with a complete cost and price-date audit still pending \todo{sum the API costs}.

Generative AI assisted author-side simulated review and manuscript editing. It did not supply new experimental results; this assistance is not formal ARR peer review.
